% LSPV-3 REPORT::KPATH1::2026-10-05T10:34:00-07:00
\documentclass[11pt]{article}
\usepackage[margin=.88in]{geometry}
\usepackage{amsmath,amssymb}
\begin{document}
\section*{LSPV-3 KPATH1: pathwise survival of the unknown SAT exponent}
\textbf{STATUS: OPEN.} P versus NP remains open. This report proves an exact accounting theorem for a declared finite reconstruction dependency DAG. It proves no unrestricted SAT lower bound and no arbitrary-algorithm acquisition theorem.

\section*{CORPUS REVIEWED}
Refetched and obeyed AGENTS.md, CODEX.md, the pnp-research skill, and root main.tex. Reviewed MAIN KFIX-SYN1, LSPV-4 KFIXGUARD1 and SEMRETGUARD1, LSPV-3 KABSORB1/RETSOURCE1/FULLSKOLEM1/SOURCESTOP1, LSPV-4 ACQ1, TASK-004, TASK-005, CUT-BOTTLENECK-AUDIT1, and repository search surfaces for HEC/CRS-PPC/CCSPC, MULTIRETURN, acquisition, PH absorption, and duplicate formulations. The newest scheduled LSPV-2 run had no verified default-branch report visible during preflight; its automation was given a write-recovery instruction without reconstructing missing scientific content.

Primary-literature preflight rechecked Williams's Circuit-SAT self-improvement theorem, the 2025 van Melkebeek--Mocelin Sdroievski targeted-generator/reconstruction paper, and alternation-trading/time-space lower-bound literature. These calibrate self-improvement and model matching but do not supply the missing ordinary-CNF return or arbitrary-static-SAT acquisition theorem.

\section*{Target and exact reconstruction guard}
Retain the quantified reconstruction interface
\[
\Psi(x,Q)=\exists \sigma\ \forall u\ \exists z\ \forall y\ R(x,Q,\sigma,u,z,y).
\]
FULLSKOLEM1 and PHABSORB1 remain semantic/fixed-level collapse statements only. KABSORB1 and KFIXGUARD1 show that implementing such collapse with positive-scale calls to a hypothetical SAT decider retains its unknown exponent \(K\).

The target here is to extend that guard from one max-recurrence to an arbitrary finite acyclic dependency graph of declared implementation branches.

\section*{PROVED/EXACT K-path theorem for a declared DAG}
Let \(G=(V,E)\) be a finite directed acyclic graph. Each node \(i\) has exponent budget \(C_i(K)\). A deterministic local stop at \(i\) has cost \(d_i\). A recursive edge \(e=(i,j)\) has nonnegative additive charge \(a_e\) and scale \(0\le\theta_e\le1\), contributing
\[
a_e+\theta_e C_j(K).
\]
A semantic leaf \(s\) reached from node \(i\) has charge
\[
a_s+\beta_s K,\qquad \beta_s\ge0.
\]
The declared budget inequalities are the maximum of all applicable deterministic stops, recursive edges, and semantic leaves.

For a directed path
\[
P=(e_1,\ldots,e_t)
\]
from root \(r\) to semantic leaf \(s\), define
\[
\Theta_P=\prod_{\ell=1}^t\theta_{e_\ell}
\]
and
\[
A_P=a_{e_1}+\theta_{e_1}a_{e_2}+\cdots+
\left(\prod_{\ell=1}^{t-1}\theta_{e_\ell}\right)a_{e_t}
+\Theta_P a_s.
\]
Then backward substitution in the finite DAG gives the least feasible root budget as the maximum over deterministic-stop paths and semantic-leaf paths. Every semantic path contributes exactly
\[
A_P+\Theta_P\beta_s K.
\]

\paragraph{Proof.}
Topologically order the DAG from leaves toward the root. At a leaf the statement is immediate. Suppose the path expansion is valid for every child \(j\) of \(i\). Substituting the child's maximum of affine path expressions into
\(a_e+\theta_e C_j(K)\) distributes the nonnegative scalar \(\theta_e\) over that maximum and prepends \(a_e\), producing exactly the path formula above. Taking the maximum over all outgoing branches and local stops completes the induction.

\paragraph{K-survival criterion.}
For this declared finite DAG, the root budget is independent of the unknown \(K\) only if every root-to-semantic-leaf path has
\[
\Theta_P\beta_s=0.
\]
If some reachable semantic leaf has \(\beta_s>0\) and every edge on one root-to-leaf path has \(\theta_e>0\), then that path contributes a strictly positive coefficient of \(K\), so the declared root budget remains unbounded as \(K\) varies.

This is an accounting theorem for the specified implementation. It does not prove that another implementation cannot avoid the semantic branch.

\section*{Cycle extension and why repeated contraction does not automatically help}
For a recursive strongly connected component with cycle product \(\Theta_\gamma<1\), RECURCLOSE1/CYCLEGAIN1 can bound repeated traversal by a finite fixed point. But if the max-recurrence also permits a positive-coefficient semantic exit before or after finitely many cycle traversals, at least one finite exit path retains a positive \(K\) coefficient. Repeatedly traversing the cycle can make a particular later-exit coefficient small, but a maximum over legal charged branches is not improved by voluntarily adding extra contractions unless the source semantics force those contractions before every semantic call.

Thus a claimed exponent erasure must prove a \emph{mandatory contraction dominator}: every path to every unresolved \(K\)-dependent call must pass through enough source-certified contraction that the resulting implementation obtains a \(K\)-free bound or a separately justified strict improvement. An optional favorable recursive branch is insufficient.

\section*{Application to the current reconstruction program}
KPATH1 converts the current instruction \emph{draw the complete dependency DAG of every term containing \(K\)} into an exact criterion. For a candidate source-specific return/acquisition construction, mark every call to the hypothetical SAT decider with its actual input-scale exponent \(\beta_s\), and mark every source-certified scale transformation by \(\theta_e\). A candidate passes the accounting gate only after every root-to-semantic-leaf path is audited.

The current vMS/SU source still does not provide an ordinary-CNF SAT-semantic return family, so no numerical return \((\eta,\theta)\) is assigned. RETSOURCE1 and RETURNTYPE1 remain controlling. The all-input reconstruction use of SAT-oracle computation cannot be treated as free PH absorption: when implemented by the assumed decider at positive scale, KABSORB1/KPATH1 retain the corresponding \(K\)-path.

\section*{Charge separation}
For a typed implementation node retain
\[
C_i\ge\max\{r_i,\chi_i,\delta_i+p_i,\text{semantic branches},\text{recursive branches}\}.
\]
Classification \(\chi_i\), setup \(\delta_i+p_i\), acquisition, generation, explicit writing, transport, rereading/recomputation, and retained space remain separate. CLASSCHARGE1 forbids adding classification to setup unless the implementation actually multiplies them. BATCHSETUP1-style sharing requires source-proved SAME\_REALIZATION grouping.

The space ledger remains independently
\[
\rho_i\ge\max\{\sigma_i,a_i,s_i,\theta_e\rho_j\}
\]
for the declared representation, with any all-retained setup floor charged only when actually required.

\section*{Hostile controls and barrier audit}
SEMRETGUARD1 is the semantic-oracle control: if the constructor learns SAT(F), constant-size return is vacuous. KPATH1 records that leak as a positive \(K\)-leaf when the answer is obtained by the assumed decider.

DEAD/COPY/DELAY/INTERP/RECOMP and padding can add or move paths and charges; they cannot justify deleting a positive \(K\)-path from a fixed declared implementation. Conversely, they demonstrate why KPATH1 is not an implementation-independent lower bound.

Horn, dual-Horn, 2-SAT, affine/Tseitin, bounded treewidth, and independent products remain easy controls: a valid universal acquisition theorem must not infer hardness merely because a path ledger contains many branches.

A maximal semantic oracle can erase the cost of semantic leaves, so KPATH1 relativizes. Algebraic-oracle access does not repair the missing SAT-semantic return. No nonrelativizing or nonalgebrizing hinge is obtained. Natural-Proofs issues are not load-bearing for this accounting theorem.

\section*{Literature transfer audit}
Williams's self-improvement theorem is a THEOREM TRANSFER only as a positive precedent that model-matched circuit-analysis problems can admit genuine bootstrapping; its construction is not transferred to ordinary CNF here. The vMS targeted-generator reconstruction is a THEOREM TRANSFER for its stated AM reconstruction semantics, not for SAT self-reduction. Alternation-trading lower bounds are model-matched terminal constraints only after an actual SAT algorithm is placed in their resource region. Max-plus/path language is merely a convenient algebraic description of the declared recurrence and supplies no complexity lower bound by itself.

\section*{DUPLICATION / CONTRADICTION AUDIT}
RECURCLOSE1 proves scalar recursive closure. KABSORB1 exposes one positive-scale SAT substitution. KFIXGUARD1 proves their one-node max composition. KPATH1 adds the finite dependency-DAG expansion and the mandatory-contraction-dominator criterion. It does not contradict KFIXGUARD1; the latter is the one-node special case. CUT-BOTTLENECK-AUDIT1 independently warns that raw event counts are not invariant obligations.

\section*{WHAT CHANGED RELATIVE TO ROOT MAIN}
Root MAIN KFIX-SYN1 asks that the smallest legal return edge expose every \(K\)-dependent branch. KPATH1 makes that requirement pathwise: a single reachable positive-scale semantic leaf through positive-scale edges suffices to retain \(K\) in the declared simulation. Therefore local optimization of one recursive edge is not load-bearing until all semantic paths are dominated by certified contractions or independently \(K\)-free algorithms.

\section*{ADOPTED / REJECTED / RETIRED PRIOR COMPONENTS}
\textbf{ADOPTED:} FULLSKOLEM1, PHABSORB1 at fixed-level semantic scope, KABSORB1, RECURCLOSE1/CYCLEGAIN1, KFIXGUARD1, SEMRETGUARD1, RETSOURCE1, RETURNTYPE1, CLASSCHARGE1, BATCHSETUP1, SAME\_REALIZATION and charge separation.

\textbf{REJECTED as sufficient:} one favorable \(\theta<1\) edge; optional repeated contraction; qualitative PH collapse; a compact evaluator treated as ordinary CNF; output contraction after learning the source decision.

\textbf{FALSIFIED/RETIRED as a generic inference:} some recursive branch contracts, therefore the unknown SAT exponent is erased. The correct declared-implementation condition is pathwise.

\section*{STALE OR OVERSTATED CLAIMS FOUND}
No verified root claim was found asserting pathwise \(K\)-erasure. The current root is conservative. The supplied historical BATCHFIX1 fallback is not used as verified evidence because current synthesis explicitly treats it as unsupported on the default branch.

\section*{STRICT STATUS}
\textbf{PROVED/EXACT:} finite-DAG path expansion and K-survival criterion for a declared implementation.

\textbf{KNOWN THEOREM:} Williams self-improvement in its circuit-analysis model; vMS reconstruction in its targeted-generator/AM model; established SAT time-space lower bounds in their stated restricted resource regions.

\textbf{CONDITIONAL:} insertion of any future source-specific return/acquisition DAG into KPATH1.

\textbf{OPEN:} exact TASK-001 source transcription and point set; source-specific ordinary-SAT return; arbitrary-algorithm acquisition; representation-invariant no-sharing; static-to-HEC compiler; nonrelativization; nonalgebrization; final contradiction; P versus NP.

\section*{SINGLE NEXT TARGET}
Take the first source-specific acquisition candidate rather than another schematic return. Produce its exact dependency DAG from original SAT input to the acquired reconstruction object. Every call to the assumed SAT decider must be labeled by actual input-scale exponent, every contraction by a source-certified scale, and every setup/classification/writing/recovery charge separately. Search first for a positive-\(K\) path that bypasses the proposed contraction. If one exists, retire that implementation before optimizing its fixed point.

\begin{thebibliography}{4}
\bibitem{Wil24} R. Ryan Williams, \emph{Self-Improvement for Circuit-Analysis Problems}, STOC 2024.
\bibitem{vMS25} D. van Melkebeek and N. Mocelin Sdroievski, \emph{Instance-Wise Hardness and Refutation versus Derandomization for Arthur--Merlin Protocols}, Computational Complexity 34:15, 2025.
\bibitem{FLMV05} L. Fortnow, R. Lipton, D. van Melkebeek, and A. Viglas, \emph{Time-Space Lower Bounds for Satisfiability}, J. ACM 52(6):835--865, 2005.
\bibitem{Wil13} R. Williams, \emph{Alternation-Trading Proofs, Linear Programming, and Lower Bounds}, ACM Trans. Comput. Theory 5(2), 2013.
\end{thebibliography}
% END_OF_VERIFIED_REPORT::LSPV-3::KPATH1
\end{document}
